% Part: first-order-logic
% Chapter: natural-deduction
% Section: proof-theoretic-notions

\documentclass[../../../include/open-logic-section]{subfiles}

\begin{document}

\iftag{FOL}
      {\olfileid{fol}{ntd}{ptn}}
      {\olfileid{pl}{ntd}{ptn}}

\olsection{నిరూపణ-సిద్ధాంత భావనలు}

\begin{editorial}
ఈ విభాగం సహజ నిగమనానికి నిరూప్యతా సంబంధం, అవైరుధ్యం యొక్క
నిర్వచనాలను సమీకరిస్తుంది.
\end{editorial}

\begin{explain}
అనేక ముఖ్యమైన అర్థపర భావనలను (చెల్లుబాటు, అనుగమనం, సంతృప్తత)
నిర్వచించినట్లే, ఇప్పుడు వాటికి అనుగుణమైన \emph{నిరూపణ-సిద్ధాంత
భావనలను} నిర్వచిస్తాం. వీటిని \tetoken{నిర్మాణాలలో}{!!{structure}s} \tetoken{వాక్యాలు}{!!{sentence}s}
సంతృప్తి చెందడాన్ని ఆశ్రయించి నిర్వచించం; కొన్ని \tetoken{వాక్యాలు}{!!{sentence}s} ఇతర
వాక్యాల నుంచి \tetoken{వ్యుత్పాద్యత}{!!{derivability}} లేదా \tetoken{అవ్యుత్పాద్యత}{!!{nonderivability}} కావడాన్ని
ఆశ్రయించి నిర్వచిస్తాం. ఈ భావనలు ఒకటే అవుతాయని కనుగొనడం ముఖ్యమైన
విజయం. అవి ఒకటే అవుతాయనేదే \emph{నిర్దుష్టత}, \emph{సంపూర్ణత}
సిద్ధాంతాల విషయం.
\end{explain}


\begin{defn}[సిద్ధాంతాలు]
అన్ని పరికల్పనలూ \tetoken{ఉపసంహరించిన}{!!{discharged}} అయిన సహజ నిగమనంలో
\tetoken{వాక్యం}{!!{sentence}}~$!A$ యొక్క ఒక \tetoken{వ్యుత్పత్తి}{!!a{derivation}} ఉంటే, $!A$ను
\emph{సిద్ధాంతం} అంటాం. $!A$ ఒక సిద్ధాంతమైతే $\Proves !A$ అని,
కాకపోతే $\Proves/ !A$ అని రాస్తాం.
\end{defn}

\begin{defn}[\tetoken{వ్యుత్పాద్యత}{!!^{derivability}}]
నిష్కర్ష~$!A$గా ఉండి, ప్రతి పరికల్పనా \tetoken{ఉపసంహరించిన}{!!{discharged}} అయినా లేదా
$\Gamma$లోనిదైనా అయ్యే ఒక \tetoken{వ్యుత్పత్తి}{!!{derivation}} ఉంటే, \tetoken{వాక్యాలు}{!!{sentence}s}
సమితి~$\Gamma$ నుంచి \tetoken{ఒక వాక్యం}{!!^a{sentence}}~$!A$ \emph{\tetoken{వ్యుత్పాదించదగిన}{!!{derivable}}}, అంటే
$\Gamma \Proves !A$. $\Gamma$ నుంచి $!A$ \tetoken{వ్యుత్పాదించదగిన}{!!{derivable}} కాకపోతే
$\Gamma \Proves/ !A$ అని రాస్తాం.
\end{defn}

\begin{defn}[అవైరుధ్యం]
\tetoken{వాక్యాలు}{!!{sentence}s} సమితి~$\Gamma$కు $\Gamma \Proves \lfalse$ అయితే,
అప్పుడు మాత్రమే అది \emph{వైరుధ్యభరితమైనది}. $\Gamma$ వైరుధ్యభరితం
కాకపోతే, అంటే $\Gamma \Proves/ \lfalse$ అయితే, దాన్ని
\emph{అవైరుధ్యమైనది} అంటాం.
\end{defn}

\begin{prop}[స్వావర్తనత్వం]
\ollabel{prop:reflexivity}
$!A \in \Gamma$ అయితే $\Gamma \Proves !A$.
\end{prop}

\begin{proof}
$!A$ పరికల్పన తనంతట తానే $!A$ యొక్క ఒక \tetoken{వ్యుత్పత్తి}{!!a{derivation}}; అందులోని ప్రతి
\tetoken{ఉపసంహరించని}{!!{undischarged}} పరికల్పన (అంటే $!A$) $\Gamma$లో ఉంటుంది.
\end{proof}


\begin{prop}[ఏకదిశత]
\ollabel{prop:monotonicity}
$\Gamma \subseteq \Delta$, $\Gamma \Proves !A$ అయితే $\Delta
\Proves !A$.
\end{prop}

\begin{proof}
$\Gamma$ నుంచి $!A$ యొక్క ఏ \tetoken{వ్యుత్పత్తి}{!!{derivation}} అయినా $\Delta$ నుంచి
$!A$ యొక్క ఒక \tetoken{వ్యుత్పత్తి}{!!a{derivation}} కూడా అవుతుంది.
\end{proof}

\begin{prop}[సంక్రామకత్వం]
\ollabel{prop:transitivity}
$\Gamma \Proves !A$, $\{!A\} \cup \Delta \Proves !B$ అయితే
$\Gamma \cup \Delta \Proves !B$.
\end{prop}

\begin{proof}
$\Gamma \Proves !A$ అయితే, $\Gamma$లోని \tetoken{ఉపసంహరించని}{!!{undischarged}}
పరికల్పనలతో $!A$ యొక్క ఒక \tetoken{వ్యుత్పత్తి}{!!a{derivation}}~$\delta_0$ ఉంటుంది.
$\{!A\} \cup \Delta \Proves !B$ అయితే, $\{!A\} \cup \Delta$లోని
\tetoken{ఉపసంహరించని}{!!{undischarged}} పరికల్పనలతో $!B$ యొక్క ఒక \tetoken{వ్యుత్పత్తి}{!!a{derivation}}~$\delta_1$
ఉంటుంది. ఇప్పుడు కిందిదాన్ని పరిశీలించండి:
\begin{prooftree}
  \AxiomC{$\Delta, \Discharge{!A}{1}$}
  \RightLabel{$\delta_1$}
  \DeduceC{$!B$}
  \DischargeRule{\Intro{\lif}}{1}
  \UnaryInfC{$!A \lif !B$}
  \AxiomC{$\Gamma$}
  \RightLabel{$\delta_0$}
  \DeduceC{$!A$}
  \RightLabel{\Elim{\lif}}
  \BinaryInfC{$!B$}
\end{prooftree}
ఇప్పుడు \tetoken{ఉపసంహరించని}{!!{undischarged}} పరికల్పనలన్నీ $\Gamma \cup \Delta$లోనివే;
కనుక ఇది $\Gamma \cup \Delta \Proves !B$ అని చూపుతుంది.
\end{proof}

$\Gamma = \{!A_1, !A_2, \ldots, !A_k\}$ ఒక పరిమిత సమితి అయితే,
$\Gamma \Proves !B$కు సరళీకృత సంజ్ఞ $!A_1,!A_2,\ldots,!A_k
\Proves !B$ను వాడవచ్చు. ముఖ్యంగా, $!A \Proves !B$ అంటే
$\{!A\} \Proves !B$.

$\Gamma \Proves !A$, $!A \Proves !B$ అయితే $\Gamma \Proves !B$
అని గమనించండి. అలాగే $!A_1, \dots, !A_n \Proves !B$ అయి, ప్రతి
$i$కి $\Gamma \Proves !A_i$ అయితే $\Gamma \Proves !B$ అని అనుగమిస్తుంది.

\begin{prop}
\ollabel{prop:incons}
కిందివి తుల్యమైనవి.
    \begin{enumerate}
      \item \( \Gamma \) వైరుధ్యభరితమైనది.
      \item ప్రతి \tetoken{వాక్యం}{!!{sentence}}~\( {!A} \)కు \( \Gamma \Proves {!A} \).
      \item ఏదో \tetoken{వాక్యం}{!!{sentence}}~\( {!A} \)కు \( \Gamma \Proves {!A} \), \( \Gamma \Proves \lnot {!A} \).
    \end{enumerate}
\end{prop}

\begin{proof}
అభ్యాసం.
\end{proof}

\tagprob{FOL}
\begin{prob}
\olref[fol][ntd][ptn]{prop:incons}ను నిరూపించండి.
\end{prob}
\tagendprob

\tagprob{notFOL}
\begin{prob}
\olref[pl][ntd][ptn]{prop:incons}ను నిరూపించండి.
\end{prob}
\tagendprob

\begin{prop}[సంహతత్వం]
  \ollabel{prop:proves-compact}
  \begin{enumerate}
  \item $\Gamma \Proves !A$ అయితే, $\Gamma_0 \Proves !A$ అయ్యే
    పరిమిత ఉపసమితి $\Gamma_0 \subseteq \Gamma$ ఒకటి ఉంటుంది.
  \item $\Gamma$లోని ప్రతి పరిమిత ఉపసమితీ అవైరుధ్యమైనదైతే,
    $\Gamma$ అవైరుధ్యమైనది.
  \end{enumerate}
\end{prop}

\begin{proof}
  \begin{enumerate}
    \item $\Gamma \Proves !A$ అయితే, $\Gamma$ నుంచి $!A$ యొక్క
      ఒక \tetoken{వ్యుత్పత్తి}{!!a{derivation}}~$\delta$ ఉంటుంది. $\delta$లోని
      \tetoken{ఉపసంహరించని}{!!{undischarged}} పరికల్పనల సమితిని $\Gamma_0$ అందాం. ప్రతి
      \tetoken{వ్యుత్పత్తి}{!!{derivation}} పరిమితం కనుక $\Gamma_0$లో పరిమిత సంఖ్యలో మాత్రమే
      \tetoken{వాక్యాలు}{!!{sentence}s} ఉండగలవు. కాబట్టి $\delta$ అనేది పరిమిత
      $\Gamma_0 \subseteq \Gamma$ నుంచి $!A$ యొక్క ఒక \tetoken{వ్యుత్పత్తి}{!!a{derivation}}.
    \item ఇది $!A \ident \lfalse$ అనే ప్రత్యేక సందర్భంలో (1)కు
      ప్రతిపక్ష నిరూపణ.
  \end{enumerate}
\end{proof}

\end{document}
